\chapter{Допуск тяжёлых стрелочных циклов горизонтальной фазы}
\label{ch:v8-horizontal-phase-heavy-arrow-cycle-admission-gate}

Вернёмся к полному бимодульному носителю
$\mathcal T_{\rm bimod}=\mathcal I_{10}\oplus\mathcal H_{10}$ и забудем
вакуумные амплитуды, сохранив только опору концевых блоков. В порядке
вершин
\begin{equation}
 (Q_L,L_L,X_L,Y_L\mid u_R,d_R,e_R,X_R,Y_R)
 \label{eq:v8-heavy-cycle-vertices}
\end{equation}
она имеет девять вершин и одиннадцать рёбер: семь инцидентностных и четыре
тяжёлых. Для ориентированной матрицы границы $\partial$ точный ранг равен
\begin{equation}
 \rank\partial=8,
 \qquad
 b_1=11-8=3.
 \label{eq:v8-heavy-cycle-betti-rank}
\end{equation}
Граф связен и действительно содержит три независимых цикла.

Инцидентностная часть сама по себе является лесом:
\begin{equation}
 \dim\ker\partial_I=7-\rank\partial_I=0.
 \label{eq:v8-heavy-cycle-incidence-forest}
\end{equation}
В порядке рёбер
\begin{multline}
 (Q_Lu_R,Q_Ld_R,L_Le_R,X_Le_R,X_LX_R,L_LY_R,Y_LY_R,\\
 X_Ld_R,Y_Le_R,L_LX_R,Q_LY_R)
 \label{eq:v8-heavy-cycle-edge-order}
\end{multline}
базис полного цикл-пространства можно выбрать столбцами матрицы
\begin{equation}
 C=\begin{pmatrix}
0&0&0\\0&0&-1\\-1&-1&1\\0&1&-1\\0&-1&0\\
1&0&-1\\-1&0&0\\0&0&1\\1&0&0\\0&1&0\\0&0&1
\end{pmatrix},
 \qquad \partial C=0,
 \qquad \rank C=3.
 \label{eq:v8-heavy-cycle-basis}
\end{equation}
Проекция последних четырёх строк $C_H$ также имеет ранг $3$. Поэтому
тяжёлые стрелки не являются спектаторами: они создают все три независимых
циклических направления.

Но это ещё не фазочувствительный объект. Единственное ребро, входящее в
$u_R$, есть $Q_Lu_R$, и оно является листовым:
\begin{equation}
 C_{Q_Lu_R,\bullet}=(0,0,0).
 \label{eq:v8-heavy-cycle-up-leaf}
\end{equation}
Следовательно, ни один замкнутый путь не проходит через сектор, который
отличает вес $4$ от общего веса $3$ остальных строк.

Тот же запрет виден непосредственно на фазовом характере. Для
$A(z)=\operatorname{diag}(z^4I_3,z^3I_7)A$ вектор весов рёбер равен
\begin{equation}
 q=(4,3,3,3,3,3,3,3,3,3,3).
 \label{eq:v8-heavy-cycle-phase-weights}
\end{equation}
Любое обычное колчанное циклическое слово чередует стрелки и сопряжённые
стрелки, а его показатель фазы задаётся циркуляцией. Точно
\begin{equation}
 qC=(0,0,0).
 \label{eq:v8-heavy-cycle-zero-charge}
\end{equation}
Таким образом, все три голономии калибровочно допустимы, но горизонтально
нейтральны. Это согласуется с общим правилом следового действия: левое
унитарное умножение сокращается в каждом замкнутом слове с $A$ и $A^*$.

\begin{theorem}[Три тяжёлых цикла без фазового заряда]
Полная бимодульная опора связна и имеет цикл-ранг $3$; инцидентностная опора
является лесом, а тяжёлые стрелки порождают все независимые циклы. Однако
ребро $Q_Lu_R$ листовое и все три циклических фазовых заряда равны нулю.
\end{theorem}

\begin{nogocriterion}[No-go обычного тяжёлого цикла]
Нельзя поднимать горизонтальную моду добавлением следа голономии любого
цикла текущего двуслойного колчана. Такой след неизбежно содержит равное
число стрелок и сопряжённых стрелок при каждой целевой вершине, поэтому
фаза сокращается. Следующий допустимый механизм должен проверить, даёт ли
Real-структура независимую ориентированную голоморфную свёртку; считать
сопряжённую стрелку новым независимым полем запрещено.
\end{nogocriterion}

\begin{protocol}
\begin{enumerate}
 \item вершин полного графа: \textbf{$9$};
 \item рёбер полного графа: \textbf{$11=7+4$};
 \item связных компонент: \textbf{$1$};
 \item цикл-ранг полной опоры: \textbf{$3$};
 \item цикл-ранг инцидентностной опоры: \textbf{$0$};
 \item ранг тяжёлой проекции цикл-пространства: \textbf{$3$};
 \item степень вершины $u_R$: \textbf{$1$};
 \item циклов через ребро $Q_Lu_R$: \textbf{$0$};
 \item горизонтальные заряды базисных циклов: \textbf{$(0,0,0)$};
 \item тяжёлые циклы поднимают горизонтальную моду: \textbf{нет};
 \item следующий гейт: \textbf{Real-ориентированная голоморфная свёртка}.
\end{enumerate}
\end{protocol}

Машинный сертификат сохранён в
\path{s2t_v8_horizontal_phase_heavy_arrow_cycle_admission_gate_results.json}.

\begin{thebibliography}{9}
\bibitem{v8-heavy-cycle-quiver-action}
C.~I.~P\'erez-S\'anchez,
\emph{The Spectral Action on Quivers}, arXiv:2401.03705.
\bibitem{v8-heavy-cycle-hashimoto}
K.~Hashimoto,
\emph{Zeta Functions of Finite Graphs and Representations of
$p$-Adic Groups}, Advanced Studies in Pure Mathematics 15 (1989),
211--280.
\end{thebibliography}